\section{Mixed-parity case analyses and two-adic lifting}
\label{sec:two-adic-lifting}

Theorem~\ref{thm:mixed-residue-classes} collects the following cases.
All residue-role permutations are included. Each row records either a
nonintegral class or a necessary condition for an integer spectrum.
\begin{center}
\begin{tabular}{p{3.2cm}p{1.6cm}p{8.8cm}}
\toprule
Residue roles & Modulus & Conclusion or necessary condition\\
\midrule
Six classes listed in the main theorem & $8$ & Nonintegral by the divisor-$32$ endpoint table below.\\
$(3,3,2)$ & $4$ & Nonintegral: endpoint and derivative requirements are incompatible.\\
$(3,0,0)$ & $4$ & Nonintegral: one is excluded by root distribution and $h_C(2)\equiv4\pmod8$.\\
$(1,3,2)$ & $4$ & $a+c\equiv b(b+2)\pmod{128}$, $b\equiv11,15\pmod{16}$ and the stated binary-cubic parity condition.\\
$(1,3,2)$, $(3,3,0)$ & $4$ & For minimum $m\geq4$, $h_C(2)=0$ and maximum $M<R(m)$.\\
\bottomrule
\end{tabular}
\end{center}

The complete identities, intermediate conditions, examples and diagnostic
boundaries are retained below. The variables identify residue roles rather
than size order; no additional refinement or search is introduced.

\subsection{Endpoint divisibility and odd-root distribution}

\begin{lemma}
\label{lem:three-odd-roots}
Let $f\in\mathbb Z[x]$ be monic of degree five with
$f(x)\equiv x^2(x+1)^3\pmod2$. If all its roots are integers,
then $32\mid f(1)$ or $32\mid f(3)$.
\end{lemma}

\begin{proof}
The binary factorization forces precisely three odd roots and two even
roots, counted with multiplicity. At least two odd roots share a residue
$r\in\{1,3\}$ modulo four. In $f(r)$, their two factors are divisible
by four and the third odd-root factor is divisible by two. The even-root
factors are odd. Therefore $32\mid f(r)$, including when $f(r)=0$.
\end{proof}

\begin{proposition}
\label{thm:mixed-parity-congruence}
Let $a,b$ be the two odd exponents and $c$ the even exponent of a positive
three-prime exponent triple. Each row in the following table, including
exchange of $a,b$, implies that its ideal intersection graph is nonintegral.
\begin{center}
\begin{tabular}{ccc|cc}
\toprule
$a\bmod8$ & $b\bmod8$ & $c\bmod8$ & $h_C(1)\bmod32$ & $h_C(3)\bmod32$\\
\midrule
1&3&2&8&16\\
1&7&2&24&16\\
3&7&2&16&8\\
3&5&6&8&16\\
3&7&6&16&24\\
5&7&6&24&16\\
\bottomrule
\end{tabular}
\end{center}
\end{proposition}

\begin{proof}
For either mixed exponent parity,
$h_C(x)\equiv x^2(x+1)^3\pmod2$. Substituting
$(a,b,c)=(8u+r_1,8v+r_2,8w+r_3)$ in the general quotient quintic gives
the last two columns: every nonconstant coefficient in $u,v,w$ is
divisible by $32$. Both entries are nonzero modulo $32$, contradicting
Lemma~\ref{lem:three-odd-roots} if all quotient roots were integers.
The resulting noninteger real complement eigenvalue has a noninteger
graph lift.
\end{proof}

For example, $(10,17,19)$ and $(14,19,21)$ meet these criteria with
gcd one and unequal $2$-adic valuations. There is no minimum or ratio
bound in the argument. The weaker divisibilities $4\mid h_C(2t)$
and $8\mid h_C(2t+1)$ hold identically for both mixed exponent parities;
they alone do not supply the contradiction. The distribution of the
three odd roots modulo four is the additional constraint.

\begin{lemma}
\label{lem:odd-root-distribution}
Let $f\in\mathbb Z[x]$ be monic of degree five with
$f(x)\equiv x^2(x+1)^3\pmod2$. Suppose
$F(y)=f(2y+1)/8$ belongs to $\mathbb Z[y]$ and
$F(y)\equiv y^3\pmod2$. If all roots of $f$ are integers,
then its three odd roots are congruent to one modulo four,
$64\mid f(1)$ and $16\mid f'(1)$.
\end{lemma}

\begin{proof}
Write the two even roots as $\eta_1,\eta_2$ and the three odd
roots as $2r_1+1,2r_2+1,2r_3+1$, with multiplicities. Then
\[
 F(y)=\prod_{j=1}^2(2y+1-\eta_j)\prod_{i=1}^3(y-r_i).
\]
The first two factors reduce to one modulo two. Hence
$\prod_i(y-r_i)\equiv y^3\pmod2$, which forces every $r_i$ to be even.
Each odd-root factor in $f(1)$ is therefore divisible by four,
giving $64\mid f(1)$. Every term in the product-rule expansion of
$f'(1)$ retains at least two such factors, giving $16\mid f'(1)$.
The argument includes coincident roots and zero factors.
\end{proof}

\begin{proposition}
\label{thm:odd-root-distribution}
Every positive three-prime exponent triple that is a permutation of
$(3,3,2)$ modulo four has a nonintegral ideal intersection graph.
\end{proposition}

\begin{proof}
Relabel the odd exponents as $a=4u+3$, $b=4v+3$ and the even
exponent as $c=4w+2$, where $u,v,w\geq0$. Expansion of the
general complement quintic gives the coefficientwise identities
\begin{align*}
 h_C(2y+1)/8&\in\mathbb Z[u,v,w,y], &
 h_C(2y+1)/8&\equiv y^3\pmod2,\\
 h_C(1)&\equiv16(u+v)\pmod{32}, &
 h_C'(1)&\equiv8(u+v+1)\pmod{16}.
\end{align*}
Its binary reduction is $x^2(x+1)^3$. If all five roots were integers,
Lemma~\ref{lem:odd-root-distribution} would give $64\mid h_C(1)$,
so the endpoint congruence would require $u+v$ even. The derivative
divisibility $16\mid h_C'(1)$ would require $u+v$ odd, a contradiction.
Thus there is a noninteger real complement eigenvalue, whose graph
lift is noninteger.
\end{proof}

This theorem contains the two rows $(3,7,2)$ and $(3,7,6)$ of
Proposition~\ref{thm:mixed-parity-congruence} and also covers odd exponents
sharing the same residue modulo eight. For example, $(10,19,27)$ and
$(14,23,39)$ have odd residues $(3,3)$ and $(7,7)$ modulo eight,
respectively. Neither theorem requires a low-root interval or an
exponent-size bound.

For the fixed necessary divisor $32$, increasing the exponent modulus
alone gives no further exclusions. On all mixed-parity triples, the
values of both $h_C(1)$ and $h_C(3)$ modulo $32$ already have period
eight in each exponent. This follows coefficientwise by substituting
the two parity representatives $(2u+1,2v,2w)$ and $(2u+1,2v+1,2w)$
and shifting one variable by $4z$; symmetry covers their permutations.
Thus the higher-modulus tables are lifts of the six classes in
Proposition~\ref{thm:mixed-parity-congruence}. The additional exclusions in
Proposition~\ref{thm:odd-root-distribution} use the distribution of the roots
and the derivative, beyond those two endpoint divisibilities.

\begin{proposition}
\label{thm:clustered-odd-roots}
Let $a,b,c$ be positive exponent residue roles with
$a\equiv1$, $b\equiv3$ and $c\equiv2\pmod4$.
If their ideal intersection graph is Laplacian integral, then
\[
 a+c\equiv b(b+2)\pmod{128},\qquad b\equiv11\text{ or }15\pmod{16}.
\]
Moreover, put $u=(a-1)/4$ and $t=(a+c-b(b+2))/128$. Then
\[
 \begin{cases}
 t+u+v\equiv1\pmod2,& b=16v+11,\\
 t+v\equiv1\pmod2,& b=16v+15.
 \end{cases}
\]
If $m=\min\{a,b,c\}\geq4$ and $M=\max\{a,b,c\}$, then also
\[
 h_C(2)=0,\qquad M<7m-16+\frac{40}{m+2}.
\]
Every violation under the stated hypotheses proves nonintegrality;
all exponent permutations are included.
\end{proposition}

\begin{proof}
Proposition~\ref{thm:odd-root-lift} in Appendix~\ref{sec:two-adic-lifting}
forces all three hypothetical integer odd quotient roots to be $b$ modulo
eight. Thus $512\mid h_C(b)$ and $64\mid h_C'(b)$.
The exact identity
\[
 h_C(b)=a^2b^2c^2\bigl(a+c-b(b+2)\bigr)
\]
has prefactor of $2$-adic valuation two, giving the sum condition modulo
$128$, including when $h_C(b)=0$. The derivative congruence and the binary
cubic calculation in Appendix~\ref{sec:two-adic-lifting} give the restriction
on $b$ and the two parity conditions.

Since $b\equiv3\pmod4$, one cannot be a quotient root under the
integer-spectrum assumption. For $m\geq4$, the positive root in $(0,3)$
from Theorem~\ref{thm:uniform-low-root} must therefore equal two.
Proposition~\ref{thm:endpoint-two-bound} gives the strict linear upper bound.
\end{proof}

The arithmetic conditions imply the earlier sum congruences
$a+c\equiv15\pmod{16}$ and $a+c+4b\equiv27\pmod{32}$.
Appendix~\ref{sec:two-adic-lifting} gives their staged derivation and the
supporting derivative thresholds. The theorem collects these arithmetic
conditions together with the endpoint-two restriction.
For example, $(29,27,754)$ and $(45,47,2258)$ have $h_C(b)=0$ and
$64\mid h_C'(b)$ but fail the final parity test. The compatible triples
$(25,27,758)$ and $(29,27,882)$ instead fail the linear upper bound.
Satisfying the arithmetic conditions alone does not establish integrality.

\begin{lemma}
\label{lem:mixed-endpoint-one}
For positive exponent triples that are permutations of $(3,0,0)$,
$(1,3,2)$ or $(3,3,0)$ modulo four, an integer complement quotient
spectrum would force every odd root to be three modulo four.
Consequently one cannot be a quotient root under this assumption.
\end{lemma}

\begin{proof}
For each displayed residue pattern, write
$(a,b,c)=(4u+r_a,4v+r_b,4w+r_c)$. Coefficientwise expansion gives
\begin{align*}
 h_C(x)&\equiv x^2(x+1)^3\pmod2,\\
 F(y)=h_C(2y+1)/8&\in\mathbb Z[u,v,w,y], &
 F(y)&\equiv(y+1)^3\pmod2.
\end{align*}
If every root were an integer, there would be three odd roots $2r_i+1$
and two even roots $\eta_j$, counted with multiplicity. In
\[
 F(y)=\prod_{j=1}^2(2y+1-\eta_j)\prod_{i=1}^3(y-r_i),
\]
the first two factors reduce to one modulo two. Thus every $r_i$ is
odd, giving the asserted residue and exclusion of one. Symmetry of the
quintic includes every permutation of the exponents.
\end{proof}

\begin{proposition}
\label{thm:one-odd-endpoints}
Every positive three-prime exponent triple that is a permutation of
$(3,0,0)$ modulo four gives a nonintegral ideal intersection graph,
without a size or ratio bound.
\end{proposition}

\begin{proof}
For residue roles $(a,b,c)=(4u+3,4v,4w)$, coefficientwise expansion
gives $h_C(2)\equiv4\pmod8$, so two is not a root.
If the minimum exponent is at least four, Theorem~\ref{thm:uniform-low-root}
provides a positive quotient root in $(0,3)$. Under the integer-spectrum
assumption it must be one or two. Lemma~\ref{lem:mixed-endpoint-one}
excludes one and the endpoint congruence excludes two, a contradiction.
The only possible smaller minimum is three, already covered by
Theorem~\ref{thm:minimum-three}. A noninteger complement root has a
noninteger graph lift.
\end{proof}

The triples $(19,36,100)$ and $(35,52,120)$ illustrate this one-odd-exponent
class beyond the minimum-through-seven and middle-through-fifteen
certificates. The result uses the uniform low-root theorem together with
the shifted root distribution, without enlarging either finite certificate.

\subsection{Derivative divisibility}

\begin{lemma}
\label{lem:derivative-thresholds}
Let $f\in\mathbb Z[x]$ be monic of degree five with two even and three
odd integer roots. If exactly $k$ of the odd roots are one modulo four,
then
\[
 2^{3+k}\mid f(1),\qquad
 2^{\max(2,k+1)}\mid f'(1),\qquad
 2^{\max(2,k)}\mid f''(1),\qquad 2\mid f'''(1).
\]
At three, the same statements hold with $k$ replaced by $3-k$.
\end{lemma}

\begin{proof}
At one, the odd-root factors supply $k$ factors divisible by four and
$3-k$ divisible by two; the even-root factors are odd. Each term of
the $j$th derivative is $j!$ times a product retaining $5-j$ root factors.
Removing the odd-root factors with the strongest guaranteed divisibility
gives the stated bounds, including the factor two from $2!$ and $3!$.
At three, the roles of the two odd residue classes are exchanged.
Repeated roots and zero evaluations obey the same divisibilities.
\end{proof}

\subsection{Odd roots modulo four and eight}

\begin{proposition}
\label{thm:odd-root-sum}
Let $a,b,c$ be positive exponents with
$a\equiv1$, $b\equiv3$ and $c\equiv2\pmod4$.
If $a+c\not\equiv15\pmod{16}$, their ideal intersection graph is
nonintegral. All permutations are included, without a size or ratio bound.
\end{proposition}

\begin{proof}
Write $(a,b,c)=(4u+1,4v+3,4w+2)$. Coefficientwise expansion gives
\[
 h_C(2y+1)/8\in\mathbb Z[u,v,w,y],\qquad
 h_C(2y+1)/8\equiv(y+1)^3\pmod2.
\]
As in Lemma~\ref{lem:odd-root-distribution}, integer quotient roots
would force all three odd roots to be three modulo four. Consequently
$64\mid h_C(3)$. A further coefficientwise identity is
\[
 h_C(3)/16\in\mathbb Z[u,v,w],\qquad
 h_C(3)/16\equiv(2v+1)(u+w+1)\pmod4.
\]
Since $2v+1$ is invertible modulo four, this forces
$u+w+1\equiv0\pmod4$, equivalently $a+c\equiv15\pmod{16}$.
Failure of that condition gives a noninteger complement root and graph lift.
\end{proof}

For example, $(14,25,27)$ and $(10,29,31)$ pass the earlier divisor-$32$
test but have $h_C(3)\equiv32\pmod{64}$, contradicting the new requirement.
The sum class $a+c\equiv15\pmod{16}$ remains compatible with this test.
Complete second- and third-derivative checks at one and three give no
additional exclusions with the root information modulo four used here.
In particular $h_C'''(1)\equiv2\pmod4$ for every mixed exponent parity.
Full valuations at a chosen exponent representative are not generally
constant throughout its residue class: $h_C(3)=-2592$ at $(1,3,6)$ but
$h_C(3)=0$ at $(9,3,6)$, although the exponent triples agree modulo eight.

\begin{proposition}
\label{thm:odd-root-lift}
Let $a,b,c$ be positive exponents with
$a\equiv1$, $b\equiv3$ and $c\equiv2\pmod4$.
If $a+c+4b\not\equiv27\pmod{32}$, their ideal intersection graph is
nonintegral. All permutations are included, without a size or ratio bound.
\end{proposition}

\begin{proof}
Suppose all five complement quotient roots are integers.
Proposition~\ref{thm:odd-root-sum} requires $a+c\equiv15\pmod{16}$.
Write
\[
 (a,b,c)=(4u+1,4v+3,16t-4u-2).
\]
Coefficientwise expansion gives
\[
 H(z)=h_C(4z+3)/64\in\mathbb Z[u,v,t,z],\qquad
 H(z)\equiv z^3+vz^2+vz+t+1\pmod2.
\]
The three odd roots are already known to be three modulo four.
Writing them as $4r_i+3$ and the two even roots as $\eta_j$ gives
\[
 H(z)=\prod_{j=1}^2(4z+3-\eta_j)\prod_{i=1}^3(z-r_i).
\]
The first two factors reduce to one modulo two, so the displayed cubic
must split into linear factors over $\mathbb F_2$.
For even $v$, it splits precisely when $t$ is odd: otherwise it is
$(z+1)(z^2+z+1)$. For odd $v$, it splits precisely when $t$ is even:
otherwise it is $z(z^2+z+1)$. The quadratic $z^2+z+1$ is irreducible.
Thus $t+v$ must be odd. Since
\[
 a+c+4b=16(t+v)+11,
\]
this forces the asserted congruence. Its failure gives a noninteger
complement root and graph lift.
\end{proof}

This excludes half the classes left by Proposition~\ref{thm:odd-root-sum}.
There are $512$ role-ordered classes modulo thirty-two; the earlier
condition leaves $128$, and the new one leaves $64$. Thus $64$ classes,
or $384$ after all permutations, are additionally excluded.
For example, the ordered triples $(13,18,19)$ and $(9,31,38)$ satisfy
the old sum condition and $64\mid h_C(3)$, but their residue-role weighted
sums are eleven modulo thirty-two. The compatible binary cubics are
$z^3$ and $(z+1)^3$, so hypothetical integer odd roots would all be
$b$ modulo eight, also requiring $512\mid h_C(b)$.
The surviving congruence classes do not assert integral spectra.

\subsection{Value, derivative and the third shifted cubic}

\begin{proof}[Proof of Proposition~\ref{thm:clustered-odd-roots}, arithmetic conditions]
The proof of Proposition~\ref{thm:odd-root-lift} forces the three hypothetical
integer odd complement roots to be $b$ modulo eight. The two remaining
roots are even. Thus $512\mid h_C(b)$ and $64\mid h_C'(b)$, since each
first-derivative term retains at least two factors divisible by eight.
The exact quintic identity
\[
 h_C(b)=a^2b^2c^2\bigl(a+c-b(b+2)\bigr)
\]
has prefactor with $2$-adic valuation exactly two. It follows that
$128\mid a+c-b(b+2)$, including when the value is zero.

Write $(a,b,c)=(4u+1,4w+3,b(b+2)+128t-a)$.
Coefficientwise expansion gives
\[
 h_C'(b)\equiv16(w^2-w+2)\pmod{64}.
\]
For $w\equiv0,1,2,3\pmod4$, the derivative residues are respectively
$32,32,0,0$. Its required divisibility forces $w\equiv2$ or $3\pmod4$,
which is the stated restriction on $b$.

For the two remaining residues, put $K(z)=h_C(8z+b)/512$.
Coefficientwise identities give $K\in\mathbb Z[u,v,t,z]$ and
\[
 K(z)\equiv
 \begin{cases}
 z^3+z^2+(1+v+u^2)z+t,&b=16v+11,\\
 z^3+z^2+(1+v)z+t,&b=16v+15
 \end{cases}\pmod2.
\]
Writing the odd roots as $b+8r_i$ and the even roots as $\eta_j$ gives
\[
 K(z)=\prod_{j=1}^2(8z+b-\eta_j)\prod_{i=1}^3(z-r_i).
\]
The even-root factors reduce to one, so the displayed binary cubic must
split completely. For $\lambda,t\in\mathbb F_2$, the cubic
$z^3+z^2+\lambda z+t$ splits precisely when $t=\lambda$: the two
possibilities are $z^2(z+1)$ and $(z+1)^3$. Otherwise it is
$z(z^2+z+1)$ or the irreducible cubic $z^3+z^2+1$.
Using $u^2\equiv u\pmod2$ gives the two stated parity conditions.
\end{proof}
